\section*{第 \ref{ch:g9:weight-vs-mass} 章——重力与质量：测量力}

\begin{solution}{exo:g9:weight-vs-mass:1}
\omterm{def:g9:weight-vs-mass:newton}{牛顿}（\unit{N}）是用\omterm{def:g9:weight-vs-mass:dynamometer}{测力计}\omterm{def:g6:measurement-in-science:measuring}{测量}的\omterm{def:g6:measurement-in-science:measuring}{单位}。参考值：
一\omterm{def:g9:weight-vs-mass:newton}{牛顿}约为一个小苹果的\omterm{def:g4:weight-and-mass:weight}{重力}；一百\omterm{def:g9:weight-vs-mass:newton}{牛顿}
约为一次有力的握手（或一个 \qty{10}{kg} 行李箱的\omterm{def:g4:weight-and-mass:weight}{重力}）。
\end{solution}

\begin{solution}{exo:g9:weight-vs-mass:2}
$P = m \times g$：\omterm{def:g4:weight-and-mass:weight}{重力}以\omterm{def:g9:weight-vs-mass:newton}{牛顿}为\omterm{def:g6:measurement-in-science:measuring}{单位}，\omterm{def:g3:measuring-mass:mass}{质量}以\omterm{def:g3:measuring-mass:units}{千克}为\omterm{def:g6:measurement-in-science:measuring}{单位}，
而 $g$——当地的\omterm{prop:g9:weight-vs-mass:formula}{重力加速度}，即每\omterm{def:g3:measuring-mass:units}{千克}物体在该处
所受拉力的\omterm{def:g9:weight-vs-mass:newton}{牛顿}数（地球上为 \qty{9.8}{N/kg}）。
\end{solution}

\begin{solution}{exo:g9:weight-vs-mass:3}
狗：$25 \times 9.8 = \qty{245}{N}$。面包：$0.7 \times 9.8
\approx \qty{6.9}{N}$。一位 \qty{55}{kg} 的读者：约
\qty{539}{N}。
\end{solution}

\begin{solution}{exo:g9:weight-vs-mass:4}
$m = 14.7 \div 9.8 = \qty{1.5}{kg}$。
\end{solution}

\begin{solution}{exo:g9:weight-vs-mass:5}
斜率即当地的 $g$——每\omterm{def:g3:measuring-mass:units}{千克}的\omterm{def:g9:weight-vs-mass:newton}{牛顿}数。在
\omterm{def:g2:moon-in-the-sky:moon}{月球}上，直线仍过原点，只是斜率降为
较平缓的 $1.6$：比例关系不变，只是\omterm{def:g4:weight-and-mass:weight}{重力}更弱。
\end{solution}

\begin{solution}{exo:g9:weight-vs-mass:6}
\omterm{def:g3:measuring-mass:mass}{质量}一列始终不变——\qty{10}{kg} 的物体
在各处都保持不变。$g$ 与\omterm{def:g4:weight-and-mass:weight}{重力}两列随
各处而变：\omterm{def:g4:weight-and-mass:weight}{重力}是当地征收的“过路费”，即 $m \times g$。
\end{solution}

\begin{solution}{exo:g9:weight-vs-mass:7}
\omterm{def:g3:levers-and-scales:scale}{天平}比较 $m_1 g$ 与 $m_2 g$：$g$ 同乘于两盘，
相互抵消——测得的\omterm{def:g3:measuring-mass:mass}{质量}处处真实。体重秤
\omterm{def:g6:measurement-in-science:measuring}{测量} $P = m g_{\text{此处}}$，却除以
$g_{\text{地球}}$ 来显示\omterm{def:g3:measuring-mass:units}{千克}数：在\omterm{def:g2:moon-in-the-sky:moon}{月球}上它显示
$m \times 1.6/9.8$——约为真实值的六分之一。
\end{solution}

\begin{solution}{exo:g9:weight-vs-mass:8}
\omterm{def:g4:weight-and-mass:weight}{重力}：\qty{19.6}{N}，竖直向下——箭头长约 \qty{3.9}{cm}
（每\omterm{def:g2:measuring-length:units}{厘米}代表 5 \omterm{def:g9:weight-vs-mass:newton}{牛顿}），从书本出发向下。
支持力：等长的 \qty{3.9}{cm} 箭头竖直向上。长度相等，
方向相反：收支平衡，书本静止。
\end{solution}

\begin{solution}{exo:g9:weight-vs-mass:9}
飞机必须承载、举起并刹停行李袋中的\emph{物质}——
即\omterm{def:g3:measuring-mass:mass}{质量}——而 \qty{23}{kg} 的物质在任何
星球上都仍是 \qty{23}{kg}。月面工作台需要一台带标定砝码的\omterm{def:g3:levers-and-scales:scale}{天平}：
唯一能经受 $g$ 变化的秤。
\end{solution}

\begin{solution}{exo:g9:weight-vs-mass:10}
$m = 111 \div 3.7 = \qty{30}{kg}$；在地球上它的\omterm{def:g4:weight-and-mass:weight}{重力}为 $30
\times 9.8 = \qty{294}{N}$。
\end{solution}

\begin{solution}{exo:g9:weight-vs-mass:11}
读数偏大：悬挂时，弹簧还要承受自身
（以及挂钩的）\omterm{def:g4:weight-and-mass:weight}{重力}，而水平放置时的零点并未计入这部分
——仪器因此多出一个恒定的偏差。将其悬挂调零，
则这一偏差在放上负载前即已扣除。
\end{solution}

\begin{solution}{exo:g9:weight-vs-mass:12}
器材：\omterm{def:g9:weight-vs-mass:dynamometer}{测力计}、一套标定砝码（$1$ 至
\qty{5}{kg}）、挂钩。逐个挂上各砝码，记录 ($m$，$P$) 数据对，
作 $P$ 关于 $m$ 的图像：各点共线并过原点，
斜率——即\omterm{def:g9:weight-vs-mass:newton}{牛顿}数增量与\omterm{def:g3:measuring-mass:units}{千克}数增量之比——为
$g \approx \qty{9.8}{N/kg}$。在\omterm{def:g2:moon-in-the-sky:moon}{月球}上：步骤相同，
同样共线，斜率变为 $1.6$——此方法
\omterm{def:g6:measurement-in-science:measuring}{测量}的是所在星球的\omterm{def:g4:weight-and-mass:weight}{重力}，变化的只是
当地的数值。
\end{solution}

\begin{solution}{pb:g9:weight-vs-mass:1}
\textbf{1.} 地球：$70 \times 9.8 = \qty{686}{N}$；\omterm{def:g2:moon-in-the-sky:moon}{月球}
\qty{112}{N}；火星 \qty{259}{N}；木星甲板
\qty{1736}{N}。
\textbf{2.} 地球上的正确\omterm{def:g3:measuring-mass:mass}{质量}为 \qty{70}{kg}。\omterm{def:g2:moon-in-the-sky:moon}{月球}：$112 \div
9.8 \approx \qty{11.4}{kg}$；火星：$259 \div 9.8 \approx
\qty{26.4}{kg}$；木星：$1736 \div 9.8 \approx
\qty{177}{kg}$——靠公式算出的，既是恭维又是诽谤。
\textbf{3.} 每个港口都是 \qty{70}{kg}：\omterm{def:g3:levers-and-scales:scale}{天平}的
比较消去了当地的 $g$。
\textbf{4.} 只有木星甲板：\qty{1736}{N} 接近
极限却仍在其下——对这位旅客而言，没有哪个港口会突破 \qty{2000}{N}
（若是体重 \qty{90}{kg} 的同事，
则在那里便会超过：$90 \times 24.8 = \qty{2232}{N}$）。
\textbf{5.} 出发时：$40 \times 9.8 = \qty{392}{N}$；
抵达时：$40 \times 3.7 = \qty{148}{N}$。
\textbf{6.} $m = 200/g$：地球 \qty{20.4}{kg}；\omterm{def:g2:moon-in-the-sky:moon}{月球}
\qty{125}{kg}；火星 \qty{54}{kg}；木星 \qty{8.1}{kg}——
工会规定的固定\omterm{def:g9:weight-vs-mass:newton}{牛顿}数，换得的\omterm{def:g3:measuring-mass:mass}{质量}却大不相同。
\textbf{7.} $m = 150 \div 1.6 = \qty{93.75}{kg}$——约
\qty{94}{kg} 的货箱，供海关核验。
\textbf{8.} $m = 294 \div 9.8 = \qty{30}{kg}$；\omterm{def:g2:moon-in-the-sky:moon}{月球}上的\omterm{def:g4:weight-and-mass:weight}{重力}
为 $30 \times 1.6 = \qty{48}{N}$。
\textbf{9.} 肌肉与\omterm{def:g3:measuring-mass:mass}{质量}回到家中时并未改变，变的只是
对手。在\omterm{def:g2:moon-in-the-sky:moon}{月球}上，举起每\omterm{def:g3:measuring-mass:units}{千克}只需
\qty{1.6}{N}；回到地球又得付 \qty{9.8}{N}。所谓增益
只是\omterm{def:g2:moon-in-the-sky:moon}{月球}的“折扣”，而非身体的成长——$P = m g$，
$m$ 忠实，$g$ 善变。
\textbf{10.} $0.2 \times 24.8 \approx \qty{5}{N}$——
它算不上铁砧之重：地球上这把打蛋器重 \qty{2}{N}，
甲板上便重 2.5 倍，对手腕颇为吃力，
但远谈不上铸造厂的活计。投诉：
部分成立。
\textbf{11.} 一\omterm{def:g9:weight-vs-mass:newton}{牛顿}的巧克力，\omterm{def:g3:measuring-mass:mass}{质量}为 $m = 1/g$：地球上约
\qty{102}{g}；\omterm{def:g2:moon-in-the-sky:moon}{月球} \qty{625}{g}；木星 \qty{40}{g}。买
在\omterm{def:g2:moon-in-the-sky:moon}{月球}上买你的\omterm{def:g9:weight-vs-mass:newton}{牛顿}吧。
\textbf{12.} 例如：“\omterm{def:g3:measuring-mass:mass}{质量}随你而行；\omterm{def:g4:weight-and-mass:weight}{重力}是
各个星球对它征收的“过路费”——带上\omterm{def:g3:measuring-mass:units}{千克}，
让各港口去收取它们的\omterm{def:g9:weight-vs-mass:newton}{牛顿}。”
\end{solution}
